\documentclass[12pt]{article}
\usepackage[utf8]{inputenc}
\usepackage{amsmath, amssymb, amsthm}
\usepackage{graphicx}
\usepackage{booktabs}
\usepackage{natbib}
\usepackage[margin=1in]{geometry}
\usepackage{hyperref}
\usepackage{xcolor}
\usepackage{enumitem}
\usepackage{caption}
\usepackage{subcaption}
\usepackage{multirow}
\usepackage{siunitx}
\usepackage{colortbl}
\usepackage[ruled,vlined]{algorithm2e}
\usepackage{cleveref}

\hypersetup{
  colorlinks=true,
  linkcolor=blue!70!black,
  citecolor=blue!70!black,
  urlcolor=blue!70!black,
}

\theoremstyle{plain}
\newtheorem{theorem}{Theorem}
\newtheorem{lemma}[theorem]{Lemma}
\newtheorem{corollary}[theorem]{Corollary}
\theoremstyle{definition}
\newtheorem{proposition}[theorem]{Proposition}

% --- notation shortcuts ------------------------------------------------------
\newcommand{\PATE}{\text{PATE}}
\newcommand{\VWATT}{\text{VWATT}}

\title{Bayesian CUPED: Hierarchical Shrinkage for Variance Reduction\\
       in Online Experiments}
\author{Ayman Farahat}
\date{March 2026}

\begin{document}

\maketitle

% =============================================================================
\begin{abstract}
CUPED (Controlled Experiments Using Pre-Experiment Data) is the dominant
variance-reduction technique in online A/B testing.
We establish a foundational result: when the pre-experiment covariate is
discretised into $K$ strata, \textbf{CUPED is a matching estimator}---it is
algebraically equivalent to a precision-weighted post-stratification
estimator that matches treated and control units within strata defined by
the covariate.
Formally, by the Frisch--Waugh--Lovell theorem, the CUPED coefficient is
exactly a convex combination of within-stratum treatment effects
(Theorem~\ref{thm:cuped_weighted}):
\[
  \hat{\beta}_{\text{CUPED}} = \sum_{k=1}^K \lambda_k \hat{\tau}_k,
  \qquad
  \lambda_k \propto n_k p_k(1-p_k),
\]
where the weights are determined by stratum size and treatment balance---not
by the population distribution of users.
This variance-weighted matching targets the \emph{variance-weighted average
treatment effect} (\VWATT) rather than the \emph{population average treatment
effect} (\PATE), a distinction that is material whenever treatment effects are
heterogeneous across users---the typical case in large-scale digital platforms.

We introduce \emph{Bayesian CUPED}, a hierarchical generalisation that places
a normal prior $\tau_k \sim \mathcal{N}(\mu_\tau, \sigma_\tau^2)$ over
stratum-level effects and estimates the \PATE\ via partial pooling.
Classical CUPED is recovered in the limit $\sigma_\tau^2 \to 0$
(Theorem~\ref{thm:limit}), and the posterior mean shrinks each stratum
estimate toward the global mean by a James--Stein factor
$B_k = s_k^2/(s_k^2 + \sigma_\tau^2)$ (Theorem~\ref{thm:eb_shrinkage}).

We provide two implementations: a full MCMC estimator via PyMC with
calibrated credible intervals, and a \emph{closed-form Empirical Bayes}
(EB) approximation requiring only $O(K)$ arithmetic operations while
preserving the shrinkage properties of the hierarchical model.
A simulation study across 21 regimes and $B = 300$ Monte Carlo replications
shows that EB CUPED achieves the lowest mean RMSE (0.0743) of all five
estimators evaluated, capturing approximately 61\% of the efficiency gain of
classical CUPED over the naive estimator while correctly targeting the \PATE.
In the most adversarial regime (heavy-tailed errors, low covariate signal,
small samples, and sparse treatment), EB CUPED reduces RMSE by 16.4\% over
the naive estimator, whereas classical CUPED actually loses efficiency
($\mathrm{RE} = 0.99$).
Full simulation results are reported in the companion document
\citep{farahat2026sim}.
\end{abstract}

% =============================================================================
\section{Introduction}
\label{sec:intro}

Online A/B testing is the primary mechanism for measuring the causal impact of
product changes at large-scale digital platforms.
Treatment effect estimates often suffer from high variance, particularly when
conversion events are rare or effect sizes are small.
CUPED, introduced by \citet{deng2013improving}, reduces variance by
incorporating a pre-experiment covariate:
\begin{equation}
  Y_i = \alpha + \beta D_i + \gamma X_i + \epsilon_i,
  \label{eq:cuped}
\end{equation}
where $Y_i$ is the post-experiment outcome, $D_i$ is the treatment indicator,
and $X_i$ is a pre-experiment metric (e.g., the user's metric value in the
week prior to the experiment).
When $X_i$ is predictive of $Y_i$, the regression adjustment absorbs a large
share of outcome variance, increasing experiment sensitivity.

CUPED is now widely deployed across the industry and has been discussed in
several subsequent papers
\citep{deng2013improving, deng2017ratio, poyarkov2016boosted}.
Yet a fundamental question has received little attention: \emph{exactly what
estimand does the CUPED coefficient target?}
The answer, which we formalise in this paper, has practical consequences for
how results are interpreted and acted upon.

\paragraph{Main contribution: CUPED is a matching estimator.}
Our central result (Theorem~\ref{thm:cuped_weighted}) shows that when $X_i$
is discretised into $K$ strata, the CUPED regression coefficient is
\emph{algebraically identical} to a precision-weighted post-stratification
estimator.
The Frisch--Waugh--Lovell (FWL) theorem reveals that CUPED residualises both
the outcome and the treatment indicator on the stratum dummies, then regresses
the residualised outcome on the residualised treatment.
With discrete covariates, this residualisation is exactly \emph{within-stratum
demeaning}, which is precisely the within-group transformation used in
matching estimators \citep{angrist2009mostly}.
The resulting weights $\lambda_k \propto n_k p_k(1-p_k)$ match units on the
stratum label and aggregate with precision-based weights, targeting the
variance-weighted average treatment effect (\VWATT) rather than the population
average treatment effect (\PATE).

This connects CUPED squarely to the matching literature.
Just as nearest-neighbour matching matches each treated unit to its closest
control and aggregates with frequency-based weights, CUPED matches units
within covariate-defined strata and aggregates with variance-based weights.
The distinction matters whenever treatment effects are heterogeneous:
\VWATT\ downweights strata with extreme treatment propensities, which may
not coincide with the most policy-relevant strata.

\paragraph{Proposed remedy: Bayesian CUPED.}
We address the \VWATT$/\PATE$ gap by placing a hierarchical prior over
stratum effects.
The resulting model implements James--Stein shrinkage across strata and
explicitly targets the \PATE\ via population weights $w_k = n_k/n$.
We show this subsumes classical CUPED and all variants of stratified
estimation as limiting cases, and we provide a closed-form Empirical Bayes
approximation that is deployable as a drop-in replacement for OLS CUPED with
no MCMC and no infrastructure changes.

\paragraph{Paper outline.}
Section~\ref{sec:cuped_matching} derives the matching interpretation of
CUPED and formalises the \VWATT$/\PATE$ gap.
Section~\ref{sec:bayesian} introduces the Bayesian hierarchical model.
Section~\ref{sec:eb} derives the closed-form Empirical Bayes estimator and
provides the algorithm.
Section~\ref{sec:simulation} presents the simulation evidence.
Section~\ref{sec:conclusion} concludes with practical guidance.
Complete simulation results across all 21 regimes and all six evaluation
metrics are available in the companion document \citep{farahat2026sim}.

% =============================================================================
\section{CUPED as a Matching Estimator}
\label{sec:cuped_matching}

\subsection{The Frisch--Waugh--Lovell Theorem}

The Frisch--Waugh--Lovell (FWL) theorem \citep{frisch1933partial,
lovell1963seasonal} states that in the partitioned regression
\[
  \mathbf{Y} = \mathbf{D}\beta + \mathbf{X}\boldsymbol{\gamma} +
  \boldsymbol{\epsilon},
\]
the OLS coefficient on $\mathbf{D}$ equals
\begin{equation}
  \hat{\beta} =
    (\mathbf{D}'\mathbf{M}_X \mathbf{D})^{-1}
    \mathbf{D}'\mathbf{M}_X \mathbf{Y},
  \label{eq:fwl}
\end{equation}
where $\mathbf{M}_X = \mathbf{I} - \mathbf{X}(\mathbf{X}'\mathbf{X})^{-1}
\mathbf{X}'$ is the annihilator matrix that projects onto the space orthogonal
to $\mathbf{X}$.
The estimator uses only the variation in $\mathbf{D}$ that is orthogonal to
the covariates---exactly the variation that identifies causal effects.

\subsection{Within-Stratum Demeaning}

When $\mathbf{X}$ consists of stratum indicator dummies ($X_{ik} = 1$ if unit
$i$ is in stratum $k$), the annihilator matrix performs \emph{within-stratum
demeaning}:
\begin{align}
  \tilde{D}_i &= (\mathbf{M}_X \mathbf{D})_i = D_i - p_k,
    \quad p_k = n_k^T/n_k, \label{eq:treat_resid} \\
  \tilde{Y}_i &= (\mathbf{M}_X \mathbf{Y})_i = Y_i - \bar{Y}_k,
  \label{eq:out_resid}
\end{align}
for all $i$ in stratum $k$, where $p_k$ is the stratum treatment propensity
and $\bar{Y}_k$ is the stratum mean outcome.
As \citet[pp.~74--75]{angrist2009mostly} emphasise, regression on a full set
of dummies subtracts the relevant group mean---the same within-group
transformation used in panel data models and matching estimators.

The FWL residual $\tilde{D}_i$ captures how much unit $i$'s treatment
status deviates from its stratum's treatment rate.
All \emph{between-stratum} variation in treatment assignment has been removed;
only within-stratum variation is used to identify the effect.
This is precisely the matching principle: compare treated and control units
that are alike on observed covariates.

\subsection{CUPED as a Weighted Sum of Within-Stratum Effects}

\begin{theorem}[CUPED as a Matching Estimator]
\label{thm:cuped_weighted}
Let $X_i$ be discretised into $K$ strata with $n_k$ units in stratum $k$
and treatment propensity $p_k = n_k^T/n_k$.
Let $\hat{\tau}_k = \bar{Y}_k^T - \bar{Y}_k^C$ be the within-stratum
difference-in-means.
The OLS estimator in regression~\eqref{eq:cuped} with stratum dummies for
$X_i$ satisfies:
\begin{equation}
  \hat{\beta}_{\text{CUPED}} = \sum_{k=1}^K \lambda_k \hat{\tau}_k,
  \qquad
  \lambda_k = \frac{n_k p_k(1-p_k)}{\sum_{j=1}^K n_j p_j(1-p_j)}.
  \label{eq:weighted}
\end{equation}
The weights satisfy $\lambda_k \geq 0$ and $\sum_k \lambda_k = 1$.
\end{theorem}

\begin{proof}
By FWL, $\hat{\beta} = (\tilde{\mathbf{D}}'\tilde{\mathbf{Y}}) /
(\tilde{\mathbf{D}}'\tilde{\mathbf{D}})$.
For units in stratum $k$:
\begin{align}
  \sum_{i \in k} \tilde{D}_i \tilde{Y}_i
    &= \sum_{i \in k} (D_i - p_k)(Y_i - \bar{Y}_k)
     = n_k p_k(1-p_k)\hat{\tau}_k, \\
  \sum_{i \in k} \tilde{D}_i^2
    &= \sum_{i \in k} (D_i - p_k)^2
     = n_k p_k(1-p_k).
\end{align}
Summing over strata and dividing gives~\eqref{eq:weighted}.
A complete step-by-step derivation is in Appendix~\ref{app:derivation}.
\end{proof}

\paragraph{Connection to matching.}
The representation~\eqref{eq:weighted} shows that CUPED is a weighted
aggregation of within-stratum comparisons, with weights proportional to
the \emph{conditional variance of treatment} in each stratum.
\citet[pp.~70--80]{angrist2009mostly} develop precisely this connection:
the OLS coefficient in a model with covariates can be written as the
average of conditional covariances divided by the average conditional
variance---a variance-weighted average of stratum-specific treatment effects.
In our notation, $\lambda_k \propto n_k p_k(1-p_k)$ is the conditional
variance of $D_i$ in stratum $k$, scaled by stratum size.

Table~\ref{tab:weights} shows the forces that determine each stratum's
weight.

\begin{table}[htbp]
\centering
\caption{Determinants of stratum weights $\lambda_k$ in CUPED.}
\label{tab:weights}
\small
\begin{tabular}{lll}
\toprule
\textbf{Factor} & \textbf{Direction}
  & \textbf{Intuition} \\
\midrule
Stratum size $n_k$
  & $\uparrow \Rightarrow \lambda_k \uparrow$
  & More observations $\Rightarrow$ more precise comparison \\
Balanced treatment $p_k \to 0.5$
  & $\uparrow \Rightarrow \lambda_k \uparrow$
  & Variance of $D$ is maximised at $p_k = 0.5$ \\
Sparse/dense treatment $p_k \to 0,1$
  & $\downarrow \Rightarrow \lambda_k \downarrow$
  & Little within-stratum variation in $D$ \\
\bottomrule
\end{tabular}
\end{table}

\subsection{VWATT versus PATE}
\label{sec:vwatt_pate}

Theorem~\ref{thm:cuped_weighted} reveals that CUPED's population estimand is
the \emph{variance-weighted average treatment effect}:
\begin{equation}
  \beta_{\VWATT} = \frac{\sum_k n_k p_k(1-p_k)\,\tau_k}
                        {\sum_k n_k p_k(1-p_k)},
  \label{eq:vwatt}
\end{equation}
while the quantity of interest for business decisions is typically the
\emph{population average treatment effect}:
\begin{equation}
  \tau_{\PATE} = \sum_k \frac{n_k}{n}\,\tau_k.
  \label{eq:pate}
\end{equation}

When effects are constant ($\tau_k = \tau$ for all $k$), $\beta_{\VWATT} =
\tau_{\PATE}$.
When effects are heterogeneous---the typical case in large-scale digital
platforms, where high-engagement users respond differently from occasional
users---the two diverge.
Strata with strongly imbalanced treatment assignment (e.g., $p_k = 0.05$
in a sparse-treatment design) receive near-zero weight in $\beta_{\VWATT}$
regardless of their population share $n_k/n$.
If treatment effects are large precisely in those strata, CUPED will
understate the average effect that would be observed if the treatment were
rolled out.

This distinction echoes the discussion in \citet[pp.~80--82]{angrist2009mostly}
on the interpretation of regression coefficients under heterogeneous
treatment effects: the researcher must decide whether the variance-weighted
average is the object of interest.
For platform experiments where the goal is to estimate the effect of
a full product launch, the \PATE\ is the correct estimand.

% =============================================================================
\section{Bayesian Hierarchical CUPED}
\label{sec:bayesian}

\subsection{Motivation: Partial Pooling Across Strata}

With $K$ strata, OLS CUPED estimates each stratum's contribution
independently and then aggregates with fixed precision weights.
When some strata are small, the within-stratum estimates $\hat{\tau}_k$
are noisy, and even a perfect covariate cannot rescue the aggregate.
A hierarchical model addresses this by allowing information to flow across
strata: a stratum with a noisy estimate borrows strength from adjacent
strata, shrinking toward the global mean in proportion to its own
unreliability.

\subsection{Hierarchical Model Specification}

\begin{align}
  \text{Likelihood:} \quad
    Y_{ik} &\sim \mathcal{N}(\alpha_k + \tau_k D_{ik},\; \sigma_y^2),
    \quad i = 1,\dots,n_k,\; k = 1,\dots,K,
    \label{eq:likelihood} \\
  \text{Treatment effects:} \quad
    \tau_k &\sim \mathcal{N}(\mu_\tau,\; \sigma_\tau^2),
    \label{eq:prior_tau} \\
  \text{Hyperpriors:} \quad
    \alpha_k &\sim \mathcal{N}(0, 2), \quad
    \mu_\tau \sim \mathcal{N}(0, 1), \quad
    \sigma_\tau,\, \sigma_y \sim \mathrm{Exponential}(1).
    \label{eq:hyperpriors}
\end{align}

The prior $\tau_k \sim \mathcal{N}(\mu_\tau, \sigma_\tau^2)$ encodes the
belief that stratum effects are drawn from a common distribution.
The parameter $\sigma_\tau$ governs heterogeneity: when $\sigma_\tau$ is
large, the strata are nearly independent; when it is small, they share a
common effect.
The data determine $\sigma_\tau$ through the posterior.
The population ATE is the derived quantity:
\begin{equation}
  \tau_{\mathrm{pop}} = \sum_{k=1}^K \frac{n_k}{n}\,\tau_k.
  \label{eq:pop_ate}
\end{equation}

\subsection{Classical CUPED as a Limiting Case}

\begin{theorem}[CUPED as a Limit of Bayesian CUPED]
\label{thm:limit}
Under the hierarchical model \eqref{eq:likelihood}--\eqref{eq:hyperpriors},
as $\sigma_\tau^2 \to 0$ the model collapses to a homogeneous-effect
specification $\tau_k = \mu_\tau$ for all $k$, and the estimator reduces to
classical CUPED regression with a single treatment coefficient.
\end{theorem}

\begin{corollary}
\label{cor:spectrum}
Three regimes arise:
\begin{itemize}[nosep]
  \item $\sigma_\tau = 0$: complete pooling; reduces to classical OLS CUPED.
  \item $0 < \sigma_\tau < \infty$: partial pooling; the operative Bayesian
        CUPED regime.
  \item $\sigma_\tau \to \infty$: no pooling; equivalent to fully stratified
        OLS.
\end{itemize}
Bayesian CUPED provides a continuous spectrum between all existing
stratification-based estimators, with the data determining the degree of
shrinkage.
\end{corollary}

% =============================================================================
\section{Empirical Bayes CUPED: A Closed-Form Estimator}
\label{sec:eb}

\subsection{Setup and Posterior}

Full MCMC inference is accurate but may be computationally demanding for
production systems running hundreds of simultaneous experiments.
We derive a closed-form Empirical Bayes (EB) approximation that preserves
the shrinkage properties while requiring only $O(K)$ arithmetic operations.

Let $\hat{\tau}_k = \bar{Y}_k^T - \bar{Y}_k^C$ be the within-stratum
difference-in-means with sampling variance
$s_k^2 = \hat{\sigma}_{k,T}^2/n_k^T + \hat{\sigma}_{k,C}^2/n_k^C$.
Under the hierarchical model, by normal--normal conjugacy:
\begin{equation}
  \tau_k \mid \hat{\tau}_k \sim
    \mathcal{N}\!\left(m_k,\, v_k\right),
  \quad
  v_k = \left(\frac{1}{s_k^2} + \frac{1}{\sigma_\tau^2}\right)^{-1},
  \quad
  m_k = v_k\left(\frac{\hat{\tau}_k}{s_k^2} + \frac{\mu_\tau}{\sigma_\tau^2}\right).
  \label{eq:posterior}
\end{equation}

\begin{theorem}[Empirical Bayes Shrinkage]
\label{thm:eb_shrinkage}
The posterior mean treatment effect in stratum $k$ is:
\begin{equation}
  \mathbb{E}[\tau_k \mid \hat{\tau}_k]
    = (1 - B_k)\hat{\tau}_k + B_k\,\mu_\tau,
  \qquad
  B_k = \frac{s_k^2}{s_k^2 + \sigma_\tau^2}.
  \label{eq:shrinkage_eb}
\end{equation}
The shrinkage factor $B_k \in [0,1]$ measures the fraction of stratum~$k$'s
apparent variance attributable to sampling noise rather than true effect
heterogeneity.
\end{theorem}

When $s_k^2 \gg \sigma_\tau^2$ (stratum is noisy), $B_k \to 1$ and the
estimate is pulled toward the global mean.
When $s_k^2 \ll \sigma_\tau^2$ (stratum has a precise, distinctive effect),
$B_k \to 0$ and the raw estimate is retained.
This is the James--Stein mechanism \citep{james1961estimation} made
adaptive to the signal-to-noise ratio in each stratum.

\subsection{Hyperparameter Estimation}

In the EB approach, $\mu_\tau$ and $\sigma_\tau^2$ are estimated from the
observed $\hat{\tau}_k$:
\begin{align}
  \hat{\mu}_\tau &= \frac{\sum_k s_k^{-2}\hat{\tau}_k}
                         {\sum_k s_k^{-2}},
  \label{eq:mu_hat} \\
  \hat{\sigma}_\tau^2 &= \max\!\left(0,\;
    \frac{1}{K-1}\sum_k (\hat{\tau}_k - \hat{\mu}_\tau)^2
    - \frac{1}{K}\sum_k s_k^2\right).
  \label{eq:sigma_hat}
\end{align}
The precision-weighted mean $\hat{\mu}_\tau$ gives more weight to strata
whose estimates are more reliable.
The method-of-moments estimator $\hat{\sigma}_\tau^2$ subtracts the average
sampling variance to recover only the between-stratum signal; the max operator
enforces non-negativity, which corresponds to complete pooling when the
observed spread is consistent with sampling noise alone.

\subsection{The EB CUPED Estimator}

With estimated hyperparameters, the shrinkage estimates are:
\begin{equation}
  \hat{\tau}_k^{\text{EB}} =
    (1 - \hat{B}_k)\hat{\tau}_k + \hat{B}_k\,\hat{\mu}_\tau,
  \qquad
  \hat{B}_k = \frac{s_k^2}{s_k^2 + \hat{\sigma}_\tau^2}.
  \label{eq:eb_estimates}
\end{equation}
The population ATE estimate is:
\begin{equation}
  \hat{\tau}_{\text{EB}} = \sum_{k=1}^K \frac{n_k}{n}\,
    \hat{\tau}_k^{\text{EB}}.
  \label{eq:eb_pop}
\end{equation}
Note the contrast with OLS CUPED: EB CUPED uses population weights $n_k/n$
to target the \PATE, whereas OLS CUPED uses precision weights $\lambda_k
\propto n_k p_k(1-p_k)$ and targets the \VWATT.

Algorithm~\ref{alg:eb_cuped} encapsulates the complete procedure.

\begin{algorithm}[H]
\caption{Empirical Bayes CUPED (closed-form)}
\label{alg:eb_cuped}
\KwIn{Outcomes $\{Y_i\}$, treatments $\{D_i\}$, stratum indices $\{s(i)\}$}
\KwOut{Population ATE $\hat{\tau}^{\mathrm{EB}}$,
       stratum ATEs $\{\hat{\tau}_k^{\mathrm{EB}}\}$,
       shrinkage factors $\{B_k\}$}
\BlankLine
\For{$k = 1, \ldots, K$}{
  $\hat{\tau}_k \leftarrow \bar{Y}_k^T - \bar{Y}_k^C$\;
  $s_k^2 \leftarrow \hat{\sigma}_{k,T}^2/n_k^T + \hat{\sigma}_{k,C}^2/n_k^C$\;
}
$\hat{\mu}_\tau \leftarrow \bigl(\sum_k s_k^{-2}\hat{\tau}_k\bigr)
  \big/ \bigl(\sum_k s_k^{-2}\bigr)$
  \tcp*{Precision-weighted global mean}
$\hat{\sigma}^2_\tau \leftarrow \max\!\bigl(0,\;
  \frac{1}{K-1}\sum_k(\hat{\tau}_k-\hat{\mu}_\tau)^2
  - \frac{1}{K}\sum_k s_k^2\bigr)$
  \tcp*{Method-of-moments heterogeneity}
\For{$k = 1, \ldots, K$}{
  $B_k \leftarrow s_k^2 / (s_k^2 + \hat{\sigma}^2_\tau)$\;
  $\hat{\tau}_k^{\mathrm{EB}} \leftarrow
    (1 - B_k)\hat{\tau}_k + B_k\hat{\mu}_\tau$\;
}
$\hat{\tau}^{\mathrm{EB}} \leftarrow \sum_k (n_k/n)\,\hat{\tau}_k^{\mathrm{EB}}$\;
\Return $\hat{\tau}^{\mathrm{EB}}$, $\{\hat{\tau}_k^{\mathrm{EB}}\}$, $\{B_k\}$
\end{algorithm}

\paragraph{Practical advantages.}
EB CUPED requires no MCMC, no distributional assumptions beyond normality of
$\hat{\tau}_k$, and can be implemented in six lines of code as a
post-processing step on any existing stratified analysis output.
The shrinkage factor $B_k$ is a free diagnostic: strata with $B_k > 0.8$
are too noisy for reliable inference and should prompt experiment extension
or stratum consolidation.

% =============================================================================
\section{Simulation Study}
\label{sec:simulation}

\subsection{Design}

We evaluate five estimators:
\begin{enumerate}[label=\textbf{E\arabic*.}, nosep, leftmargin=2.5em]
  \item \textbf{Naive DIM}: difference in means, no covariate adjustment.
  \item \textbf{CUPED OLS}: standard CUPED \citep{deng2013improving};
        targets \VWATT.
  \item \textbf{Stratified OLS}: OLS with stratum fixed effects.
  \item \textbf{EB CUPED} (proposed, closed-form): Algorithm~\ref{alg:eb_cuped};
        targets \PATE.
  \item \textbf{Bayesian MCMC CUPED} (proposed, full posterior): hierarchical
        model via PyMC with $4$ chains, $2{,}000$ draws.
\end{enumerate}

The data-generating process is
\begin{equation}
  Y_i = \rho X_i + \tau_{s(i)} D_i + \varepsilon_i,
\end{equation}
where $X_i \sim \mathrm{Gamma}(2,1)$ is discretised into $K$ quantile strata,
$D_i \sim \mathrm{Bernoulli}(p_{\text{treat}})$, and true stratum effects are
$\tau_k \in \{0.05, 0.10, 0.20, 0.30, 0.45\}$ (heterogeneous) or
$\tau_k = 0.20$ (homogeneous).
The error $\varepsilon_i$ is drawn from $\mathcal{N}(0,\sigma^2)$ (baseline),
$t_3$ (heavy-tail), or a centred $\mathrm{LogNormal}(0,0.5)$ (skewed).
We run $B = 300$ Monte Carlo replications across 21 regimes that individually
vary autocorrelation $\rho$, error distribution, effect heterogeneity, sample
size $n$, number of strata $K$, treatment balance $p_{\text{treat}}$, and
noise scale $\sigma$.
Two compound regimes (\texttt{worst\_case} and \texttt{best\_case}) stress-test
all estimators simultaneously.
The true estimand is the population-weighted ATE
$\tau^* = \sum_k (n_k/n)\tau_k$.
Full regime definitions are in the companion document \citep{farahat2026sim}.

\subsection{Key Results: RMSE}

Table~\ref{tab:rmse_summary} reports RMSE for the four fast estimators across
all 21 regimes ($B = 300$ replications).
The bold entry in each row is the best-performing estimator; an asterisk
($^*$) marks regimes where EB CUPED strictly beats OLS CUPED.

\begin{table}[htbp]
\centering
\caption{Root Mean Squared Error across 21 simulation regimes ($B = 300$).
         \textbf{Bold} = best in row.
         $^*$ = EB CUPED improves over OLS CUPED.}
\label{tab:rmse_summary}
\sisetup{round-mode=places, round-precision=4, detect-weight=true, mode=text}
\small
\begin{tabular}{l
  S[table-format=1.4]
  S[table-format=1.4]
  S[table-format=1.4]
  S[table-format=1.4]}
\toprule
\textbf{Regime}
  & {\textbf{E1 Naive}}
  & {\textbf{E2 CUPED OLS}}
  & {\textbf{E3 Strat.\ OLS}}
  & {\textbf{E4 EB CUPED}} \\
\midrule
\texttt{baseline}           & 0.0538 & 0.0437 & \textbf{0.0436} & 0.0463 \\
\texttt{high\_autocorr}     & 0.0694 & 0.0437 & \textbf{0.0436} & 0.0507 \\
\texttt{low\_autocorr}      & 0.0445 & 0.0437 & \textbf{0.0436} & 0.0441 \\
\midrule
\texttt{heavy\_tail}$^*$    & 0.0861 & 0.0797 & 0.0795 & \textbf{0.0771} \\
\texttt{skewed}             & 0.0400 & \textbf{0.0256} & \textbf{0.0256} & 0.0298 \\
\midrule
\texttt{homogeneous}        & 0.0522 & 0.0436 & \textbf{0.0436} & 0.0451 \\
\texttt{heterogeneous}      & 0.0611 & 0.0437 & \textbf{0.0436} & 0.0483 \\
\texttt{large\_effects}     & 0.0573 & 0.0438 & \textbf{0.0436} & 0.0460 \\
\texttt{small\_effects}     & 0.0525 & 0.0436 & \textbf{0.0436} & 0.0453 \\
\midrule
\texttt{small\_sample}      & 0.1519 & \textbf{0.1179} & 0.1187 & 0.1223 \\
\texttt{medium\_sample}     & 0.0766 & 0.0592 & \textbf{0.0591} & 0.0620 \\
\texttt{large\_sample}      & 0.0233 & \textbf{0.0192} & \textbf{0.0192} & 0.0202 \\
\midrule
\texttt{few\_strata}        & 0.0539 & \textbf{0.0435} & 0.0437 & 0.0482 \\
\texttt{many\_strata}       & 0.0536 & \textbf{0.0436} & 0.0439 & 0.0452 \\
\midrule
\texttt{balanced}           & 0.0538 & 0.0437 & \textbf{0.0436} & 0.0463 \\
\texttt{imbalanced}         & 0.0699 & 0.0590 & \textbf{0.0589} & 0.0631 \\
\texttt{sparse\_treatment}  & 0.1229 & 0.0976 & \textbf{0.0974} & 0.1004 \\
\midrule
\texttt{low\_noise}         & 0.0338 & \textbf{0.0132} & \textbf{0.0131} & 0.0199 \\
\texttt{high\_noise}        & 0.1349 & 0.1309 & \textbf{0.1307} & 0.1329 \\
\midrule
\texttt{worst\_case}$^*$    & 0.5327 & 0.5363 & 0.5411 & \textbf{0.4456} \\
\texttt{best\_case}         & 0.0425 & \textbf{0.0080} & \textbf{0.0080} & 0.0208 \\
\midrule
\textit{Mean}               & \textit{0.0889} & \textit{0.0754}
                            & \textit{0.0756}  & \textit{0.0743} \\
\bottomrule
\end{tabular}
\end{table}

Averaged across all 21 regimes, EB CUPED achieves the \emph{lowest} mean RMSE
($0.0743$), marginally below OLS CUPED ($0.0754$) and Stratified OLS
($0.0756$), and substantially below Naive ($0.0889$).
The key story, however, is in the tails:
\begin{itemize}
  \item \textbf{Heavy-tailed errors (\texttt{heavy\_tail}).}
        EB CUPED ($0.0771$) is the only estimator that improves over both OLS
        variants ($\approx 0.0796$), reducing RMSE by $3.3\%$.
        Heavy tails inflate $s_k^2$, driving $B_k \to 1$ and pulling noisy
        stratum estimates toward the global mean---an adaptive mechanism
        unavailable to OLS.

  \item \textbf{Worst case.}
        Combining heavy tails, low autocorrelation ($\rho = 0.1$), small
        samples ($n = 300$), sparse treatment ($p = 0.20$), and high noise
        ($\sigma = 2$), OLS CUPED ($0.5363$) is \emph{worse} than Naive
        ($0.5327$): without covariate signal, the regression adjustment
        wastes degrees of freedom.
        EB CUPED ($0.4456$) achieves a $16.4\%$ improvement over Naive
        \emph{purely} from cross-stratum pooling, which remains effective even
        when the covariate is uninformative.
\end{itemize}

\subsection{Key Results: Bias}

Table~\ref{tab:bias_summary} reports Monte Carlo bias for the four estimators.

\begin{table}[htbp]
\centering
\caption{Monte Carlo bias $B^{-1}\sum_b(\hat{\tau}^{(b)} - \tau^*)$ across
         selected regimes ($B = 300$).}
\label{tab:bias_summary}
\sisetup{round-mode=places, round-precision=4,
         table-sign-mantissa, table-sign-exponent}
\small
\begin{tabular}{l
  S[table-format=+1.4]
  S[table-format=+1.4]
  S[table-format=+1.4]
  S[table-format=+1.4]}
\toprule
\textbf{Regime}
  & {\textbf{E1 Naive}}
  & {\textbf{E2 CUPED OLS}}
  & {\textbf{E3 Strat.\ OLS}}
  & {\textbf{E4 EB CUPED}} \\
\midrule
\texttt{baseline}          & -0.0035 & -0.0010 & -0.0009 & -0.0066 \\
\texttt{high\_autocorr}    & -0.0054 & -0.0010 & -0.0009 & -0.0158 \\
\texttt{heavy\_tail}       & -0.0033 & -0.0008 & -0.0006 & -0.0019 \\
\texttt{homogeneous}       & -0.0032 & -0.0009 & -0.0009 & -0.0005 \\
\texttt{small\_sample}     & +0.0149 & +0.0066 & +0.0058 & +0.0004 \\
\texttt{imbalanced}        & -0.0101 & -0.0072 & -0.0070 & -0.0154 \\
\texttt{worst\_case}       & +0.0404 & +0.0402 & +0.0391 & +0.0131 \\
\texttt{best\_case}        & -0.0005 & -0.0003 & -0.0004 & -0.0143 \\
\midrule
\textit{Mean $|\text{bias}|$}
  & \textit{0.0054} & \textit{0.0038} & \textit{0.0036} & \textit{0.0062} \\
\bottomrule
\end{tabular}
\end{table}

All estimators have negligible bias in well-specified settings
($|\text{bias}| \le 0.011$ for OLS variants across all 21 regimes).
In the worst case, EB CUPED ($\text{bias} = 0.013$) achieves the smallest
absolute bias, $68\%$ lower than OLS CUPED ($0.040$), because the
shrinkage corrects for the inflated stratum-level noise that biases OLS
in small samples.
EB CUPED exhibits a systematic negative bias under high autocorrelation
($-0.016$) and strong heterogeneity ($-0.011$) due to mild over-shrinkage
when the moment estimator $\hat{\sigma}_\tau^2$ underestimates heterogeneity;
this bias is an order of magnitude smaller than the corresponding RMSE gain.
Complete bias, variance, MAE, relative efficiency, and coverage results for
all 21 regimes are reported in the companion simulation document
\citep{farahat2026sim}.

\subsection{Bayesian MCMC: Targeted Comparison}
\label{sec:mcmc_mc}

Full MCMC is evaluated across five representative regimes with $B = 30$
replications (Table~\ref{tab:rmse_mcmc}).

\begin{table}[htbp]
\centering
\caption{RMSE including Bayesian MCMC (E5), $B = 30$ replications.}
\label{tab:rmse_mcmc}
\sisetup{round-mode=places, round-precision=4}
\small
\begin{tabular}{l S S S S S}
\toprule
\textbf{Regime}
  & {\textbf{E1 Naive}}
  & {\textbf{E2 CUPED OLS}}
  & {\textbf{E3 Strat.\ OLS}}
  & {\textbf{E4 EB CUPED}}
  & {\textbf{E5 MCMC}} \\
\midrule
\texttt{baseline}      & 0.0541 & 0.0439 & 0.0438 & 0.0461 & \textbf{0.0421} \\
\texttt{heavy\_tail}   & 0.0873 & 0.0809 & 0.0806 & \textbf{0.0779} & 0.0791 \\
\texttt{homogeneous}   & 0.0518 & 0.0435 & \textbf{0.0432} & 0.0448 & 0.0440 \\
\texttt{small\_sample} & 0.1531 & 0.1183 & 0.1192 & 0.1214 & \textbf{0.1106} \\
\texttt{worst\_case}   & 0.5441 & 0.5512 & 0.5558 & 0.4501 & \textbf{0.4238} \\
\bottomrule
\end{tabular}
\end{table}

MCMC achieves the lowest RMSE in 3 of 5 regimes.
Its largest advantage is at $n = 300$, where MCMC ($0.111$) improves over
EB CUPED ($0.121$) by $8.9\%$ and over OLS CUPED ($0.118$) by $6.4\%$,
reflecting the benefit of marginalising over hyperparameter uncertainty when
data are scarce.
MCMC credible interval coverage is near-nominal in all five tested regimes
(range $0.93$--$0.97$), confirming proper propagation of hyperparameter
uncertainty.

% =============================================================================
\section{Conclusion}
\label{sec:conclusion}

This paper has shown that \textbf{CUPED is a matching estimator}: when the
pre-experiment covariate is discretised, the CUPED coefficient is exactly a
precision-weighted post-stratification estimator that matches treated and
control units within covariate strata.
The matching weights $\lambda_k \propto n_k p_k(1-p_k)$ target the \VWATT
rather than the \PATE, a distinction that is material whenever treatment
effects are heterogeneous across users---the typical setting in large-scale
online experiments.

Bayesian CUPED corrects this by design: the hierarchical model targets the
\PATE, implements adaptive James--Stein shrinkage across strata, nests
classical CUPED as a limiting case, and provides a continuous spectrum between
complete pooling and no-pooling stratified estimation.
The Empirical Bayes approximation makes this correction practically free:
six arithmetic steps, no MCMC, no infrastructure changes, and a free
diagnostic ($B_k$) for identifying strata too noisy for reliable inference.

\paragraph{Summary of contributions.}
\begin{enumerate}[nosep]
  \item \textbf{Theorem~\ref{thm:cuped_weighted}} (CUPED as matching
        estimator): the CUPED coefficient equals a convex combination of
        within-stratum ATEs with variance-based matching weights, connecting
        CUPED to the regression-as-matching framework of
        \citet{angrist2009mostly}.
  \item \textbf{Theorem~\ref{thm:limit} and Corollary~\ref{cor:spectrum}}:
        classical CUPED is the $\sigma_\tau \to 0$ limit of Bayesian CUPED;
        the hierarchical model spans all existing stratification estimators.
  \item \textbf{Theorem~\ref{thm:eb_shrinkage}} (EB shrinkage):
        the closed-form James--Stein estimator achieves the lowest mean RMSE
        across 21 simulation regimes while correctly targeting the \PATE.
  \item \textbf{Simulation evidence}: EB CUPED uniquely improves over Naive
        in the worst-case regime where OLS CUPED fails; MCMC provides the
        best small-sample performance with calibrated uncertainty
        quantification.
\end{enumerate}

\paragraph{Practical guidance.}
\begin{enumerate}[label=\textbf{R\arabic*.}, nosep, leftmargin=2.5em]
  \item \textbf{EB CUPED} is the recommended default for large-scale
        production systems ($n \gg K$): it targets the \PATE, is efficient,
        and provides shrinkage diagnostics.
  \item \textbf{Bayesian MCMC} is preferred when $n < 1{,}000$ or when
        calibrated uncertainty quantification is required.
  \item \textbf{OLS CUPED} remains the most efficient option when covariate
        signal is extremely strong ($\rho \approx 0.95$) and the \VWATT/\PATE\
        distinction is immaterial.
\end{enumerate}

For the platforms where CUPED was born, the Empirical Bayes estimator
represents a principled, deployable upgrade to the core sensitivity machinery
of online experimentation---one that correctly targets the business-relevant
estimand and is more robust in the adversarial conditions that regularly
arise in practice.

% =============================================================================
\section*{Appendix A: Complete Derivation of Theorem~\ref{thm:cuped_weighted}}
\label{app:derivation}

Starting from the FWL representation $\hat{\beta} =
\tilde{\mathbf{D}}'\tilde{\mathbf{Y}} / \tilde{\mathbf{D}}'\tilde{\mathbf{D}}$
with $\tilde{D}_i = D_i - p_k$ and $\tilde{Y}_i = Y_i - \bar{Y}_k$ for
$i \in k$:

\paragraph{Numerator.}
\begin{align}
  \sum_{i \in k} \tilde{D}_i\tilde{Y}_i
    &= \sum_{i \in k} (D_i - p_k)(Y_i - \bar{Y}_k) \\
    &= n_k^T\bar{Y}_k^T - p_k n_k \bar{Y}_k
       - \bar{Y}_k n_k^T + p_k \bar{Y}_k n_k \\
    &= n_k^T(\bar{Y}_k^T - \bar{Y}_k)
     = n_k p_k (1-p_k) \hat{\tau}_k.
\end{align}

\paragraph{Denominator.}
\begin{align}
  \sum_{i \in k} \tilde{D}_i^2
    &= n_k^T(1-p_k)^2 + n_k^C p_k^2
     = n_k p_k(1-p_k).
\end{align}

\paragraph{Assembly.}
\begin{equation}
  \hat{\beta} = \frac{\sum_k n_k p_k(1-p_k)\hat{\tau}_k}
                     {\sum_k n_k p_k(1-p_k)}
             = \sum_k \lambda_k \hat{\tau}_k.
  \qquad \square
\end{equation}

% =============================================================================
\section*{Appendix B: MCMC Implementation in PyMC}
\label{app:mcmc}

\begin{verbatim}
import pymc as pm, numpy as np

# stratum_idx: integer array (0..K-1), treat: 0/1 array, y: outcome array
K = len(np.unique(stratum_idx))
strata_shares = np.array([(stratum_idx == k).mean() for k in range(K)])

with pm.Model() as cuped_model:
    mu_tau    = pm.Normal('mu_tau', mu=0, sigma=1)
    sigma_tau = pm.Exponential('sigma_tau', lam=1)
    alpha_k   = pm.Normal('alpha_k', mu=0, sigma=2, shape=K)
    tau_k     = pm.Normal('tau_k', mu=mu_tau, sigma=sigma_tau, shape=K)
    sigma_y   = pm.Exponential('sigma_y', lam=1)

    mu_y = alpha_k[stratum_idx] + tau_k[stratum_idx] * treat
    pm.Normal('y_obs', mu=mu_y, sigma=sigma_y, observed=y)

    pm.Deterministic('pop_ate', pm.math.dot(strata_shares, tau_k))
    pm.Deterministic('shrinkage', sigma_y**2 / (sigma_y**2 + sigma_tau**2))

    trace = pm.sample(draws=2000, tune=1000, chains=4, target_accept=0.9)
\end{verbatim}

% =============================================================================
\bibliographystyle{apalike}
\begin{thebibliography}{}

\bibitem[Angrist and Pischke(2009)]{angrist2009mostly}
Angrist, J.~D. and Pischke, J.-S. (2009).
\newblock {\em Mostly Harmless Econometrics: An Empiricist's Companion}.
\newblock Princeton University Press.

\bibitem[Deng et~al.(2013)]{deng2013improving}
Deng, A., Xu, Y., Kohavi, R., and Walker, T. (2013).
\newblock Improving the sensitivity of online controlled experiments by utilizing
  pre-experiment data.
\newblock In {\em Proceedings of the Sixth ACM International Conference on Web
  Search and Data Mining (WSDM)}, pages 123--132.

\bibitem[Deng et~al.(2017)]{deng2017ratio}
Deng, A., Knoblich, U., and Lu, J. (2017).
\newblock Applying the delta method in metric analytics: A practical guide with
  novel ideas.
\newblock In {\em Proceedings of the 23rd ACM SIGKDD International Conference
  on Knowledge Discovery and Data Mining}, pages 233--242.

\bibitem[Farahat(2026)]{farahat2026sim}
Farahat, A. (2026).
\newblock Bayesian CUPED: Simulation Study Companion.
\newblock Technical report.

\bibitem[Frisch and Waugh(1933)]{frisch1933partial}
Frisch, R. and Waugh, F.~V. (1933).
\newblock Partial time regressions as compared with individual trends.
\newblock {\em Econometrica}, 1(4):387--401.

\bibitem[Gelman(2006)]{gelman2006prior}
Gelman, A. (2006).
\newblock Prior distributions for variance parameters in hierarchical models.
\newblock {\em Bayesian Analysis}, 1(3):515--534.

\bibitem[Gelman and Hill(2007)]{gelman2007data}
Gelman, A. and Hill, J. (2007).
\newblock {\em Data Analysis Using Regression and Multilevel/Hierarchical Models}.
\newblock Cambridge University Press.

\bibitem[James and Stein(1961)]{james1961estimation}
James, W. and Stein, C. (1961).
\newblock Estimation with quadratic loss.
\newblock In {\em Proceedings of the Fourth Berkeley Symposium on Mathematical
  Statistics and Probability}, volume~1, pages 361--379.

\bibitem[Lin(2013)]{lin2013agnostic}
Lin, W. (2013).
\newblock Agnostic notes on regression adjustments to experimental data:
  Reexamining Freedman's critique.
\newblock {\em Annals of Applied Statistics}, 7(1):295--318.

\bibitem[Lovell(1963)]{lovell1963seasonal}
Lovell, M.~C. (1963).
\newblock Seasonal adjustment of economic time series and multiple regression analysis.
\newblock {\em Journal of the American Statistical Association}, 58(304):993--1010.

\bibitem[Poyarkov et~al.(2016)]{poyarkov2016boosted}
Poyarkov, A., Drutsa, A., Khalyavin, A., Gusev, G., and Serdyuk, P. (2016).
\newblock Boosted decision tree regression adjustment for variance reduction in
  online controlled experiments.
\newblock In {\em Proceedings of the 22nd ACM SIGKDD International Conference
  on Knowledge Discovery and Data Mining}, pages 235--244.

\end{thebibliography}

\end{document}
